\section{Model}\label{sec:model}

A stylized model captures what the set-aside does and, as important, what it does not do. $N$ potential bidders compete for a single indivisible contract. Each bidder $i$ has type $k_i \in \{\text{SME}, \neg\}$ and draws a private cost of supply $c_i$ from a type-specific distribution $F_c^k$, with $F_c^{\text{SME}}$ and $F_c^{\neg}$ differing in ways the data are allowed to determine.

The auction format is the identifying payload. BEC implements procurement through a Preg\~ao---a descending-clock electronic English-reverse auction, distinct both from first-price sealed-bid designs (Convite, in the smaller-value Brazilian segment) and from the Dutch descending-price format more commonly studied in the auction literature. The clock opens at the buyer's reference price $p^{\mathrm{ref}}$, firms post downward-revisable offers until only one firm remains, and that firm supplies the contract at the clearing price. What makes this format analytically convenient is a standard result of \citet{haile2003}: under independent private values, the weakly dominant strategy is to exit the auction when the running clock price reaches one's own cost. Drop-out prices of losers therefore \emph{point-identify} their costs, without requiring the researcher to invert a bid function or assume a parametric form for the bidding strategy. The winner, who never drops out before the clock stops, leaves only an upper bound $c_{\text{win}} \le c_{(2)}$ (the cost of the second-to-last firm to exit)---and that upper bound is exactly what the Turnbull robustness check in Section~\ref{sec:robustness} uses to deliver a non-parametric MLE of $F_c^k$. This point-identification under Preg\~ao is the reason a structural analysis is available here where the cross-sectional first-price-sealed-bid settings of \citet{marion2007}, \citet{krasnokutskaya2011}, and \citet{atheyseira2011} have relied on cross-equation constraints to disentangle channels.

The model relies on three substantive restrictions, stated formally and defended empirically in the data.

\begin{assumption}[Independent Private Values, asymmetric across types]\label{a:ipv}
Costs $c_i$ are drawn independently across bidders. Conditional on type $k$, $c_i \sim F_c^k$ with density $f_c^k$ on support $[\underline{c}, \overline{c}] \subset \mathbb{R}_+$. The SME and non-SME distributions are not constrained to coincide.
\end{assumption}

Assumption~\ref{a:ipv} permits SMEs and non-SMEs to differ in cost distributions---a realistic feature given that non-SMEs are typically the large distributors of the same generic pharmaceutical---without requiring one to stochastically dominate the other.

\begin{assumption}[Multiplicative auction-level heterogeneity]\label{a:uh}
The normalized log-bid decomposes as $y_{it} = \log(b_{it} / p_t^{\mathrm{ref}}) = a_t + e_{it}$, where $a_t$ is an auction-level scale common to all bidders in auction $t$, and $e_{it}$ is bidder-specific and independent across bidders within $t$. Both $a_t$ and $e_{it}$ are mean-zero with finite variances $\sigma_a^{2,k}$ and $\sigma_e^{2,k}$ that may vary by stratum.
\end{assumption}

Assumption~\ref{a:uh} is the \citet{krasnokutskaya2011} restriction that separates a common scale from the bidder-specific private value. Multiplicativity is the natural form when the auction-level shock affects all bidders proportionally to their cost realizations---standard in procurement settings where item complexity, ANVISA registration burden, or contract specifications scale costs uniformly across firms. It is weaker than the fully non-parametric \citet{kotlarski1967} deconvolution, which would additionally require type-symmetry within $e_{it}$; Section~\ref{sec:robustness} reports a Turnbull NPMLE that dispenses with Assumption~\ref{a:uh} altogether and corroborates the baseline numbers within 16 percent of the total effect across classes (with the largest gap in pharmaceutical Turnbull).

\begin{assumption}[Set-aside, with equilibrium SME selection]\label{a:setaside}
Under the open regime both SMEs and non-SMEs enter with strictly positive probability. Under the SME-only regime, only SMEs are admissible. Within a short window around the cutoff, the underlying technology of supply is stable, but the equilibrium distribution of active SMEs under the restricted regime may differ from its pre-regime counterpart because the policy changes which SMEs find entry profitable. Write these distributions $F_c^{\text{SME},\tau}$ for $\tau \in \{\text{Pre},\text{Post}\}$.
\end{assumption}

Assumption~\ref{a:setaside} is the structural counterpart of the parallel-trends assumption in reduced-form DiD, but with equilibrium selection made explicit. In the \citet{atheyseira2011} reading adopted here, the policy changes the admissible set and therefore the composition of the SME bidders who actually enter. The post-policy SME cost distribution is thus an equilibrium object, not evidence of a primitive technological break. This interpretation fits the diagnostic evidence: the primitive-invariance KS test on non-SME $F_c$ (Table~\ref{tab:v3_primitive_invariance}) rejects strict invariance in the pharma Preg\~ao stratum ($D = 0.0722$) and in the non-pharma Convite stratum ($D = 0.0514$), and fails to reject in the pharma Convite ($D = 0.0324$) and non-pharma Preg\~ao ($D = 0.0486$) strata. The operational threshold $D < 0.05$ corresponds to less than five percentage points of CDF movement at the Kolmogorov sup norm. The cross-modality check between Convite GPV and Preg\~ao drop-out recoveries in Table~\ref{tab:v3_cross_modality} provides an additional consistency test.

\paragraph{Strict invariance as a limiting case.} A stronger reading would impose $F_c^{\text{SME,Post}} = F_c^{\text{SME,Pre}}$ and attribute the entire policy response to bidder counts alone. I treat that as a useful benchmark rather than the main specification. Section~\ref{sec:robustness} reports it explicitly as a strict-primitive-invariance limit: it preserves the intensive-margin headline, but it changes the welfare comparison in the thin, heterogeneous pharma market because it shuts down equilibrium selection in who enters.

\paragraph{What the KS rejection of strict invariance implies for the modeling choice.} The KS tests in Tables~\ref{tab:v3_primitive_invariance} and~\ref{tab:v3_uh_invariance} reject strict invariance of non-SME $F_c$ in three of the four Preg\~ao strata after UH cleaning, with KS distances of 0.0613 (non-pharma Preg\~ao) and 0.1407 (pharma Preg\~ao). This rejection is not a problem for the structural exercise; it is the empirical foundation for adopting equilibrium selection (Assumption~\ref{a:setaside}) as the main specification rather than strict invariance. A model that imposed strict invariance and then ran the BNE counterfactual on data that visibly violate it would be misspecified by construction. The main-specification approach taken here---to treat $F_c^{\text{SME,Post}}$ as an equilibrium object generated by the protected regime and to compare results against the strict-invariance bound in robustness---is exactly the response the KS evidence calls for. The rejection of strict invariance is therefore evidence in favor of the modeling choice, not against it.

\paragraph{Equilibrium and counterfactuals.} The auction format is interpreted as ascending-clock IPV throughout: under Assumption~\ref{a:ipv}, exit-at-cost is the weakly dominant strategy, the surviving bidder pays the second-lowest exit, and the expected transaction price equals $\mathbb{E}[c_{(2)}]$. Revenue equivalence \citep{vickrey1961, maskinriley2000} extends the same expected-price object to FPSB and to other revenue-equivalent formats, but the operative interpretation in this paper is the ascending clock that BEC's Preg\~ao implements; \citet{krasnokutskaya2011} provides the analog argument used here. The counterfactuals therefore reduce to Monte Carlo simulation over draws from the estimated $F_c^k$: draw a stratum-specific count of bidders from each type's Poisson arrival process, draw costs, sort, and read off the two lowest order statistics. The lowest, $c_{(1)}$, is the production cost; the second-lowest, $c_{(2)}$, is the expected transaction price. The exercise is conducted \emph{under observed equilibrium entry}: post-policy participation rates and the post-policy SME cost distribution are taken from data, not derived from a free-entry condition jointly with the cost primitive. The model is therefore not a fully solved nested entry-bidding equilibrium of the kind developed in \citet{krasnokutskayaseim2011}, where entry costs are structurally identified from a free-entry indifference condition that jointly disciplines the bidding stage. Here, entry costs are reported as zero-profit calibration margins (Section~\ref{sec:results}); the decomposition that follows is a counterfactual-simulation accounting under the modeling assumptions stated rather than a separation of independently identified structural channels.

\paragraph{Entry.} The Bayes-Nash counterfactual takes equilibrium arrival rates as primitives rather than solving for them explicitly. Observed Pre-period counts $(N^{\text{SME}}_{\text{Pre}}, N^{\neg}_{\text{Pre}})$ are the status-quo arrival rates under the open regime; Post-period counts $(N^{\text{SME}}_{\text{Post}}, 0)$ are the equilibrium arrival under the restricted regime. This is a reduced-form treatment of entry in the \citet{atheyseira2011} spirit; the main specification combines the observed Post counts with the observed Post SME bidder distribution. The implied zero-profit entry costs per bidder, reported separately (Table~\ref{tab:v3_entry_cost}), are in the R\$0.11--R\$2.46 range. These are per-auction zero-profit margins at the entry margin, not the fixed cost of preparing a BEC documentation packet (which is amortized over many auctions and is not directly recovered by this calibration); they are a consistency check on the assumed Poisson primitives rather than a structural estimate.\footnote{Independence-across-auctions is assumed at the level of cost draws conditional on type and stratum. Within-auction correlation enters the model through the auction-level scale $a_t$ (Assumption~\ref{a:uh}); the cluster bootstrap of Section~\ref{sec:results} resamples at the auction level to preserve this correlation in the standard errors.}

\paragraph{The decomposition.} Three counterfactual scenarios anchor the exercise. $S_1$ fixes the regime as open and the pool as pre-period; it is the baseline. $S_2$ imposes SME-only while holding the SME pool at its pre-policy size: the non-SME type is removed and the SME pool is held at its observed pre-policy arrival rate. $S_3$ imposes SME-only at the observed post-period SME pool size, using the observed post-policy arrival rate and the post-policy SME cost distribution. The simulated total effect on the second-order statistic decomposes arithmetically as
\begin{equation}
\underbrace{p_{S_3} - p_{S_1}}_{\text{simulated total effect}}
\;=\; \underbrace{(p_{S_2} - p_{S_1})}_{\text{within-auction component (sheltered bidding)}}
   \,+\, \underbrace{(p_{S_3} - p_{S_2})}_{\text{participation-margin component}}.
\end{equation}
Both components are properties of the simulation given $F_c^k$ and the observed entry pattern; they sum to the simulated total effect by construction. The within-auction component can exceed the realized total effect when the participation-margin component partially offsets it, which is what the data deliver in non-pharmaceutical markets.

\paragraph{Why ``sheltered bidding''.} I refer to the within-auction component as \emph{sheltered bidding}: a descriptive label for the change in the simulated $\mathbb{E}[c_{(2)}]$ when non-SMEs are removed and the SME pool is held at its pre-policy size. The label captures the equilibrium price formation under reduced competition, not any active behavioral adjustment by SMEs---under ascending-clock IPV, the dominant strategy is exit-at-cost regardless of pool composition, and the within-auction component is mechanically a property of which order statistic is realized at each pool composition. The share of the simulated total effect attributable to the within-auction component depends on the cost-distribution estimator and on whether one imposes strict invariance on the post-policy SME pool; reading the share as a tightly pinned point estimate would be misleading. The pharmaceutical decomposition is reported alongside as a more model-sensitive boundary case.

\paragraph{Scope of the structural exercise.} The model deliberately stops short of five extensions. (i) Entry is taken as observed Poisson arrival rates, not solved jointly with bidding under a free-entry condition. (ii) The implied $\kappa^k$ are zero-profit margins at the entry margin, not structurally identified fixed costs. (iii) Procurement quantity is treated as inelastic at observed values, with the implied bound on the welfare cost discussed in Section~\ref{sec:discussion}. (iv) The 10 percent preference V3 is computed under Vickrey-equivalent translation; the FPSB equivalent under \citet{maskinriley2000} adjustment is bounded above by the type-specific cost-distribution dispersion in each stratum, with the welfare ranking V3 $\succ$ V0 preserved at the upper bound (Section~\ref{sec:fpsb_bound}). (v) Asymmetric IPV is assumed; symmetric IPV is consistent with the data only at very low quantile dispersion across types and is not formally tested. Each is a margin a fully solved structural model would price; each is bracketed below or in robustness rather than imposed as the central exercise.
